% Referencias que no están en los reportes del repositorio. Verificadas el 15 de octubre de 2026.

\bibitem{jensen2017}
N. Jensen y C. Barrett, ``Agricultural Index Insurance for Development,'' \emph{Applied Economic Perspectives and Policy}, vol.~39, no.~2, pp.~199--219, 2017.

\bibitem{agrocostos}
Fideicomisos Instituidos en Relación con la Agricultura, ``Agrocostos.'' [En línea]. Disponible: \url{https://www.fira.gob.mx/Nd/Agrocostos.jsp}. Consultado el 15 de octubre de 2026.

% Documentos del equipo. Enlaces a la rama main del repositorio.
\bibitem{doc_contexto}
J.~A. Pomares Valdés, V.~J. López Martínez y L.~E. Ramiro Luna, ``El Reto AgroCebada 2026 en contexto. Lo que sabemos,
lo que suponemos y lo que falta comprobar,'' Benemérita Universidad Autónoma de Puebla, Puebla, México, Rep. técnico,
oct. 2026. [En línea]. Disponible: \url{https://github.com/alan-pv/AgroCebada_Hackaton/blob/main/AgroCebada_contexto.pdf}

\bibitem{doc_modelo}
J.~A. Pomares Valdés, V.~J. López Martínez y L.~E. Ramiro Luna, ``Rendimiento de cebada de temporal con régimen de
etiquetas y rezago espacial. Procedimiento, arquitectura y evaluación del modelo,'' Benemérita Universidad Autónoma de
Puebla, Puebla, México, Rep. técnico, oct. 2026. [En línea]. Disponible:
\url{https://github.com/alan-pv/AgroCebada_Hackaton/blob/main/docs/reporte_modelo/main.pdf}

\bibitem{doc_generalizable}
J.~A. Pomares Valdés, V.~J. López Martínez y L.~E. Ramiro Luna, ``¿Dónde rinde más la cebada de temporal? Un modelo
municipal con datos públicos para Hidalgo, México, Puebla y Tlaxcala,'' Benemérita Universidad Autónoma de Puebla,
Puebla, México, Rep. técnico, oct. 2026. [En línea]. Disponible:
\url{https://github.com/alan-pv/AgroCebada_Hackaton/blob/main/docs/reporte_generalizable/modelo_generalizable.pdf}

\bibitem{doc_regimenes}
J.~A. Pomares Valdés, V.~J. López Martínez y L.~E. Ramiro Luna, ``¿Qué separa a los dos regímenes de rendimiento?
Descarte de causas ambientales y contraste con las cifras oficiales del SIAP,'' Benemérita Universidad Autónoma de
Puebla, Puebla, México, Rep. técnico, oct. 2026. [En línea]. Disponible:
\url{https://github.com/alan-pv/AgroCebada_Hackaton/blob/main/docs/reporte_regimenes/main.pdf}

\bibitem{repo}
J.~A. Pomares Valdés, V.~J. López Martínez y L.~E. Ramiro Luna, ``AgroCebada. Código, datos públicos y reportes del Reto AgroCebada FIRA 2026,'' GitHub, 2026.
[En línea]. Disponible: \url{https://github.com/alan-pv/AgroCebada_Hackaton}
